\section{逆矩阵与 LU 分解：撤销变换并复用消元}
\label{sec:03-Linear-Algebra/01-Linear-Systems-and-Matrices/03}

\subsection{当你需要解几十遍同一个矩阵}

假想你在设计一个电路。电阻网络的拓扑不变，但你要测试许多组不同的输入电压，看输出电流怎么变。每次都是一样的系数矩阵 \(A\)，变的是右端 \(b\)。你用 Gaussian 消元老老实实算了前三个右端之后，发现了一个尴尬的事实：每次消元的步骤几乎一模一样——约掉第一列、约掉第二列……唯一变化的是增广列。

你可能会想：有没有办法把消元的过程存起来？就像做菜时提前备好料，客人来了只需要下锅翻炒两下？

有。两条路都可以走：

\begin{itemize}
\tightlist
\item 找逆矩阵 \(A^{-1}\)：如果知道``撤销 \(A\) 的变换''，直接 \(x=A^{-1}b\) 就完事；
\item 做 LU 分解 \(A=LU\)：把 \(A\) 拆成两个三角矩阵，之后每个新右端只需要两次三角求解。
\end{itemize}

两条路都是把``与右端无关的计算''提前做完。但用途不同，代价也不同。我们一个一个来看。

\subsection{什么样的变换可以撤销}

你穿上外套（变换 \(A\)），走到室外。回来时想把外套脱掉（变换 \(A^{-1}\)），回到原来的状态。能脱的前提是：外套确实穿上了，而且脱的方法唯一。

对方阵 \(A\in\R^{n\times n}\)，如果存在矩阵 \(B\) 使得

\[
AB=BA=I,
\]

就说 \(A\) 可逆，\(B\) 是 \(A\) 的逆矩阵，记作 \(A^{-1}\)。这样的 \(B\) 如果存在，一定是唯一的：若 \(B_1,B_2\) 都满足条件，则

\[
B_1=B_1I=B_1(AB_2)=(B_1A)B_2=IB_2=B_2.
\]

一旦知道了 \(A^{-1}\)，\(Ax=b\) 就有了形式上的``一步求解''：

\[
x=A^{-1}b.
\]

但``形式上的''这三个字有讲究。\(A^{-1}b\) 这个写法很漂亮，实际计算时却不一定先求 \(A^{-1}\) 再乘 \(b\)——后者比直接解方程通常更慢，也更不稳定。逆矩阵是理论分析的好工具，不一定是数值计算的好方案。

\subsection{不可逆意味着什么：信息丢在了路上}

不是所有方阵都可逆。来看看一个反例：

\[
A=\begin{pmatrix}1&2\\2&4\end{pmatrix}.
\]

假设 \(Ax=b\) 有解。两列 \((1,2)^{\trans}\) 和 \((2,4)^{\trans}\) 落在同一条直线上——列空间只是一条线，不是整个平面。现在考虑齐次方程 \(Ax=0\)：

\[
\begin{pmatrix}1&2\\2&4\end{pmatrix}
\begin{pmatrix}x\\y\end{pmatrix}
=\begin{pmatrix}0\\0\end{pmatrix}
\implies
\begin{cases}
x+2y=0,\\
2x+4y=0,
\end{cases}
\]

两个方程说的是同一件事：\(x=-2y\)。所以任何形如 \(t(-2,1)^{\trans}\) 的向量都满足 \(Ax=0\)。现在假设有一个 ``逆矩阵'' \(B\) 存在，那么

\[
BA=I
\implies
BA\begin{pmatrix}-2\\1\end{pmatrix}=B\cdot0=0,
\]

但同时

\[
I\begin{pmatrix}-2\\1\end{pmatrix}=\begin{pmatrix}-2\\1\end{pmatrix}\neq0.
\]

矛盾。所以 \(A\) 不可逆。它的几何原因是：变换 \(A\) 把多个不同的输入压到了同一个输出——\(x\) 和 \(x+t(-2,1)^{\trans}\) 经过 \(A\) 后完全一样。信息丢失了，你无法从输出唯一地恢复输入。

对方阵，以下说法描述的是同一件事（等价条件）：

\begin{itemize}
\tightlist
\item \(A\) 可逆；
\item Gaussian 消元时每一行、每一列都有主元（秩 = \(n\)）；
\item \(Ax=0\) 只有零解；
\item 对每一个 \(b\in\R^n\)，\(Ax=b\) 有唯一解；
\item \(A\) 的 \(n\) 个列可以唯一地表示 \(\R^n\) 中的任意向量。
\end{itemize}

这些等价的严格证明放到向量空间单元，那时有了维数和秩-零化度定理，整个逻辑链条会非常简洁。

\subsection{逆不是逐元素取倒数}

你可能会想：矩阵求逆，不就是把每个元素取倒数吗？试一下：

\[
A=\begin{pmatrix}2&1\\1&1\end{pmatrix}.
\]

如果逐元素取倒数，你会得到 \((\begin{smallmatrix}1/2&1\\1&1\end{smallmatrix})\)，但这乘上 \(A\) 显然不是单位阵。实际的逆矩阵是

\[
A^{-1}=\begin{pmatrix}1&-1\\-1&2\end{pmatrix}.
\]

验一下：\(AA^{-1}=\begin{pmatrix}2&1\\1&1\end{pmatrix}\begin{pmatrix}1&-1\\-1&2\end{pmatrix}=\begin{pmatrix}1&0\\0&1\end{pmatrix}=I\)。

为什么逆矩阵的元素看起来和原矩阵没有简单的倒数关系？因为求逆依赖的是各列之间的整体关系——你要找的是一组系数，使得各列的线性组合构成标准基向量。每个元素的位置和关系都参与了计算，不能孤立地对待。

复合变换的逆会反转顺序，这个直觉很自然：先穿外套再穿雨衣，脱的时候先脱雨衣再脱外套。

\[
(AB)^{-1}=B^{-1}A^{-1}.
\]

验证：\((AB)(B^{-1}A^{-1})=A(BB^{-1})A^{-1}=AIA^{-1}=AA^{-1}=I\)。反过来的乘积同理。

\subsection{换个角度看消元：每一步都是一个矩阵}

在进入 LU 分解之前，先掌握一个关键工具：\emph{初等矩阵}。

对单位矩阵做一次初等行操作，得到的矩阵就是初等矩阵。它左乘任何尺寸兼容的矩阵时，会在那个矩阵上复制同一个行操作。

例如，要在 \(2\times2\) 矩阵上执行``第二行减去第一行的 2 倍''，对应的初等矩阵是

\[
E=
\begin{pmatrix}
1&0\\
-2&1
\end{pmatrix}.
\]

把它乘到 \(A=\begin{pmatrix}2&1\\4&3\end{pmatrix}\) 上：

\[
EA=
\begin{pmatrix}
1&0\\
-2&1
\end{pmatrix}
\begin{pmatrix}2&1\\4&3\end{pmatrix}
=
\begin{pmatrix}
2&1\\
0&1
\end{pmatrix}.
\]

第一行原样保留（因为 \(E\) 的第一行是 \((1,0)\)，取出 \(A\) 的第一行），第二行变成了 \(A\) 的第二行减去 \(A\) 的第一行的 2 倍——正是 Gaussian 消元的一步。消元结果记作 \(U\)（上三角，upper triangular）。

\(E\) 是可逆的——撤销``减去第一行的 2 倍''只需``加上第一行的 2 倍''：

\[
E^{-1}=
\begin{pmatrix}
1&0\\
2&1
\end{pmatrix}.
\]

于是

\[
A=E^{-1}U.
\]

\subsection{消元过程怎样浓缩为 \(A=LU\)}

把 \(E^{-1}\) 记作 \(L\)（下三角，lower triangular）。上面的例子：

\[
\underbrace{\begin{pmatrix}2&1\\4&3\end{pmatrix}}_{A}
=
\underbrace{\begin{pmatrix}1&0\\2&1\end{pmatrix}}_{L}
\underbrace{\begin{pmatrix}2&1\\0&1\end{pmatrix}}_{U}.
\]

\(U\) 是消元的终点——上三角矩阵，主对角线下方全是零。\(L\) 记录的是\emph{为了到达 \(U\)，每一步用了前一行的多少倍}。注意 \(L\) 中的元素是\emph{乘子}（multiplier）：第二行第一列的位置是 2，意思是``消元时第二行减掉了第一行的 2 倍''。而消元矩阵 \(E\) 中同一位置是 \(-2\)。符号刚好相反——这是个非常容易犯错的细节。记住：\(L\) 存的是乘子本身，\(E\) 存的是乘子的相反数。

推广到更大的矩阵。假设消元过程不需要交换行，每一步都从对角位置取主元，消去下方的元素。那么整串消元可以写成一串初等矩阵的乘积：

\[
E_k\cdots E_2E_1 A = U,
\]

其中每个 \(E_i\) 消去第 \(i\) 列主元下方的元素。反过来：

\[
A = (E_1^{-1}E_2^{-1}\cdots E_k^{-1})U = LU.
\]

多个下三角矩阵的逆相乘，结果仍是下三角。而且 \(L\) 的对角线全为 1（因为每个 \(E_i^{-1}\) 的对角线都是 1），非对角元素就是对应位置的消元乘子。后面的消元不会干扰前面已经置好的乘子位置，所以你可以边消元边往 \(L\) 里填数。

如果消元过程中需要交换行，多一个置换矩阵 \(P\)：\(PA=LU\)，稍后单说。

\subsection{分解一次，喂给许多右端}

有了 \(A=LU\)，解 \(Ax=b\) 变成两步：

\[
LUx=b.
\]

先令 \(y=Ux\)，解

\[
Ly=b
\]

得到 \(y\)；再解

\[
Ux=y
\]

得到 \(x\)。前者是下三角系统，用\emph{前代}（forward substitution）；后者是上三角系统，用\emph{回代}（backward substitution）。

沿用刚才的矩阵，设 \(b=(1,5)^{\trans}\)：

\[
\begin{pmatrix}1&0\\2&1\end{pmatrix}
\begin{pmatrix}y_1\\ y_2\end{pmatrix}
=
\begin{pmatrix}1\\5\end{pmatrix}.
\]

第一行直接给出 \(y_1=1\)。第二行：\(2y_1+y_2=5\) 即 \(2\cdot1+y_2=5\)，得 \(y_2=3\)。

再解：

\[
\begin{pmatrix}2&1\\0&1\end{pmatrix}
\begin{pmatrix}x_1\\ x_2\end{pmatrix}
=
\begin{pmatrix}1\\3\end{pmatrix}.
\]

第二行：\(x_2=3\)。第一行：\(2x_1+x_2=1\) 即 \(2x_1+3=1\)，得 \(x_1=-1\)。

现在换一个右端 \(b'=(3,7)^{\trans}\)。\(L\) 和 \(U\) 完全不用重算。只需重做前代和回代：

\(y_1=3\)，\(2\cdot3+y_2=7\implies y_2=1\)。然后 \(x_2=1\)，\(2x_1+1=3\implies x_1=1\)。几行算完。

计算量对比：对稠密 \(n\times n\) 矩阵，LU 分解本身大约需要 \(\frac{2}{3}n^3\) 次浮点运算。分解完成后，每个新右端的前代和回代各只需约 \(n^2\) 次运算。当 \(n=1000\) 时，分解约 6.7 亿次运算，之后每个右端只需约 200 万次——差了三百多倍。这就是为什么工程师愿意花时间算 LU 分解：一次性投资，后续求解几乎免费。

你可能会问：求 \(A^{-1}\) 然后乘 \(b\)，不也是 \(n^2\) 吗？对，但求 \(A^{-1}\) 本身需要大约 \(2n^3\) 次运算（约是 LU 分解的三倍），而且逆矩阵通常比 \(L\) 和 \(U\) 稠密得多，存储和乘法都更贵。除非你需要反复用同一个逆矩阵做别的事（比如计算敏感度），否则 LU 分解 + 三角求解始终是解方程组的首选方案。

\subsection{行交换不能忽略：置换矩阵 \(P\)}

前面的 LU 分解假设消元过程中从未需要交换行。但看看这个矩阵：

\[
A=\begin{pmatrix}0&1\\1&1\end{pmatrix}.
\]

左上角是 0，无法做主元。交换两行后，新矩阵的左上角变成 1，消元可以正常进行。这个矩阵明明可逆（两列不共线），却不能做无交换的 LU 分解。

把行交换记录为一个\emph{置换矩阵} \(P\)。\(P\) 的每一行、每一列恰好有一个 1，其余为 0。左乘 \(P\) 的效果就是重新排列行。交换两行的 \(P\) 是把单位矩阵的对应两行交换得到的。

实际的分解写成

\[
PA=LU.
\]

解 \(Ax=b\) 时，两端左乘 \(P\)：

\[
PAx=Pb\quad\Longrightarrow\quad LUx=Pb.
\]

仍然是前代和回代，只不过右端先被 \(P\) 重新排列了一下。

不只是零主元需要交换。在浮点计算中，即使当前主元非零，如果它比同列其他候选小很多（比如 \(10^{-16}\) 对 1.5），用它消元会放大舍入误差。部分选主元策略会主动把绝对值最大的候选换上来，哪怕原来的主元在精确算术中完全可用。\(P\) 记录的交换中，相当一部分不是为了救零主元，而是为了数值稳定。

\subsection{无置换 LU 何时存在}

如果想确保 \(A=LU\) 存在（无需 \(P\)），也就是消元每一步的对角主元都非零，一个充分必要条件是：所有\emph{顺序主子式}（左上角 \(1\times1, 2\times2, \ldots, n\times n\) 子矩阵的行列式）都非零。行列式在后面单元正式引入，这里只需知道它衡量一个方阵是否``退化''。

对称正定矩阵（比如协方差矩阵、刚度矩阵）满足更强的性质：不但所有顺序主子式非零，而且可以分解为 \(LL^{\trans}\)（Cholesky 分解）——储存量比普通 LU 省一半，计算量也省一半。

当矩阵非常稀疏时（大多数元素为零），选择行列顺序不只是为了数值稳定——错误的顺序可能在分解过程中产生大量\emph{填充}（fill-in），把本来稀疏的 \(L\) 和 \(U\) 填得密密麻麻，存储和计算成本急剧上升。这时置换同时服务于稳定性和稀疏性两个目标，选主元的策略远比课本上的简单交换复杂。

\subsection{可逆不等于靠谱：一个诚实的区分}

一个矩阵在精确数学中完全可逆，但在有限精度的计算机上可能表现得几乎奇异。经典例子是 Hilbert 矩阵：

\[
H_3=\begin{pmatrix}
1 & 1/2 & 1/3\\
1/2 & 1/3 & 1/4\\
1/3 & 1/4 & 1/5
\end{pmatrix}.
\]

\(H_3\) 精确可逆，但它的各列非常接近彼此的线性组合。如果你对右端 \(b\) 做百万分之一的扰动，解 \(x\) 可能改变好几个数量级。你算出残差 \(\norm{Ax-b}_2\) 很小，就以为解很准——其实残差小只说明你找到了一个使 \(Ax\) 接近 \(b\) 的 \(x\)，不说明这个 \(x\) 接近真实解。

这里有两条线必须分清楚：

\begin{itemize}
\tightlist
\item \emph{可逆性}（invertibility）：回答``精确解是否存在且唯一''——这是代数结构问题；
\item \emph{条件性}（conditioning）：回答``解对数据误差敏不敏感''——这是数值分析问题。
\end{itemize}

前者看秩是否满，后者看条件数（最大奇异值与最小奇异值之比）是否大。一个矩阵可以满秩但病态，也可以良态但奇异（比如秩亏 1 的矩阵，在特定右端下解对微小扰动并不敏感）。两条线各自独立，碰巧相交时最棘手——既满秩又良态的矩阵，才是数值计算最喜欢的对象。条件数的正式定义放到数值分析单元。

延伸阅读：MIT 18.06：Factorization into \(A=LU\)。
